\documentclass[12pt]{article}
\usepackage[left=1.5in, right=1.25in, top=1.25in, bottom=1.25in]{geometry}
\usepackage{setspace}
\usepackage{booktabs}
\usepackage{times}
\doublespacing

\begin{document}

\begin{center}
{\large \textbf{JazzyFS: Revised Thesis Approach}}\\[0.2in]
{\normalsize Michelle Gurovith\\
University of California, Santa Cruz\\
June 2026}
\end{center}

\section*{Revised Positioning}

Prior work including the Linux kernel readahead, GFPP, and SPP demonstrates
that adaptive prefetching is not a new problem. The Linux kernel already detects
sequential versus random access patterns and adjusts prefetching accordingly.
Advanced research systems such as GFPP use deep graph neural networks to
produce confidence scores internally, while SPP uses path confidence to throttle
prefetching at the microarchitectural level.

However, none of these systems expose their prediction uncertainty as a
first-class observable signal at the filesystem level. SPP keeps confidence
internal to the microarchitecture. GFPP requires training data and operates
as a black-box model. No existing filesystem-level prefetcher logs
per-decision confidence values as a user-accessible signal, exposes a
tunable confidence threshold at runtime without kernel modification, or
uses sonification specifically as a prefetching observability interface.

Traditional filesystem prefetchers already approximate confidence implicitly:
the Linux readahead window grows when sequential reads succeed and shrinks when
they fail, but this reaction is delayed by several read cycles and the signal
driving it is never named or exposed. JazzyFS formalizes this implicit signal as
an explicit, per-decision confidence value and uses it as a \textbf{direct
control mechanism}. JazzyFS is therefore repositioned not as a new prefetching
algorithm, but as a \textbf{transparent, training-free, cross-platform platform}
that replaces implicit heuristic control with explicit confidence-driven control,
reducing unnecessary prefetching and improving responsiveness to workload changes
without kernel modification.

\section*{Novel Contributions}

Research confirms three contributions that are novel at the filesystem level:

\begin{enumerate}

    \item \textbf{Sonification as a Human-in-the-Loop Prefetching Interface.}
    JazzyFS is the first system to use intentional sonification as a
    prefetching observability interface and to lay the groundwork for a
    human-in-the-loop feedback loop: the operator hears confidence drop,
    intervenes, and the system adjusts future prefetching behavior. No
    existing filesystem or prefetcher --- kernel or research --- has
    implemented this closed-loop approach. Prior acoustic research concerns
    unintentional hardware noise for security side-channels, not deliberate
    operator-facing observability by design.

    \item \textbf{Confidence Decay Rate as a Phase Change Metric.}
    Existing prefetchers track confidence values but do not measure how fast
    confidence drops when a workload changes. JazzyFS exposes this decay rate
    as a first-class metric: a fast drop indicates an abrupt phase change
    while a slow drop indicates gradual drift. No filesystem prefetcher
    currently distinguishes between these two cases. This distinction has
    direct system impact: detecting an abrupt drop enables immediate
    suppression of prefetching, reducing cache pollution and wasted I/O
    during sudden transitions. Linux readahead's implicit heuristic requires
    multiple failed reads before the window shrinks; JazzyFS's explicit
    signal enables the same decision in a single cycle. This metric is
    derivable directly from JazzyFS's existing per-decision logs without
    additional experiments.

    \item \textbf{Confidence Trajectory Classification Without Training.}
    Every existing adaptive prefetcher uses confidence to answer the same
    question: what block or file comes next? JazzyFS uses the temporal shape
    of the confidence curve --- how it rises, falls, or oscillates over time
    --- to answer a fundamentally different question: what kind of workload
    is currently running? A fast rise signals sequential access; a value near
    zero signals random access; oscillation signals mixed access. This
    inversion has direct system impact: a system that classifies its workload
    as random at file-open time can suppress prefetching immediately rather
    than waiting for confidence to degrade read by read. Existing prefetchers
    react to evidence of failure; trajectory classification enables proactive
    control before failures accumulate. This classification has no precedent
    at the filesystem level and is derivable directly from JazzyFS's existing
    per-decision logs without new experiments.

\end{enumerate}

\section*{Revised Research Questions}

\begin{enumerate}

    \item \textbf{FUSE Overhead vs.\ Adaptive Logic Cost.}
    What is the overhead cost of making prefetching decisions explicit and
    observable in user-space FUSE, and is that overhead dominated by the
    FUSE interposition layer or by the adaptive logic itself?

    \item \textbf{Confidence Threshold Sensitivity.}
    How does the confidence threshold affect prefetching effectiveness across
    different workload types, and what threshold minimizes overhead while
    maintaining prefetch benefit on sequential workloads?

    \item \textbf{Confidence Decay Rate and Phase Transition Speed.}
    How quickly does a sliding-window confidence scorer detect workload phase
    transitions, and does the rate of confidence decay distinguish abrupt
    phase changes from gradual workload drift?

    \item \textbf{Confidence Trajectory Classification.}
    Does the temporal shape of the confidence curve --- how it rises, falls,
    or oscillates --- reliably classify the workload type currently running,
    and does this training-free classification match the ground-truth workload
    labels assigned during experiment construction?

    \item \textbf{Sonification as a Human-in-the-Loop Interface.}
    Does intentional sonification of prefetching confidence produce an
    auditory signal operators can act on, and does it lay the groundwork
    for a closed-loop system where an operator hears confidence drop,
    intervenes, and the prefetcher adjusts --- something no existing
    filesystem or prefetcher has implemented?

\end{enumerate}

\section*{Experimental Changes}

\subsection*{Experiments Removed}

\begin{tabular}{lp{10cm}}
\toprule
Removed Experiment & Reason \\
\midrule
Prefetch depth sweep (1, 2, 4, 8) &
Results show depth greater than one provides no benefit and only increases
overhead. The conclusion is always depth~1 is optimal, which does not
advance the thesis contribution. \\[0.2in]
Baseline OBL as primary comparison &
The kernel already performs unconditional readahead. OBL inside FUSE
replicates existing kernel behavior and does not represent a novel
comparison point given the revised framing. \\
\bottomrule
\end{tabular}

\subsection*{Experiments Added}

\begin{tabular}{lp{9cm}}
\toprule
New Experiment & Description \\
\midrule
Confidence threshold sweep
(0.3, 0.5, 0.7, 0.9) &
Run all seven workloads at each threshold value using the existing
experiment harness. Measure execution time, prefetch rate, and overhead
at each threshold. Directly supports the tunable confidence contribution.
Replaces the depth sweep with the same experimental structure. \\[0.2in]
Confidence decay rate analysis &
Using existing decision logs from the phase\_change workload, measure
how many reads it takes for confidence to drop below the threshold after
the pattern switches from sequential to random. Compute the decay rate
and classify it as abrupt or gradual. No new experiment runs required ---
existing logs contain this data. \\[0.2in]
Confidence trajectory classification &
Using existing decision logs, extract the confidence curve for each
workload run and classify its shape as rising (sequential), near-zero
(random), or oscillating (mixed). Compare classifications against
ground-truth workload labels. No new experiment runs required ---
existing logs contain this data. \\
\bottomrule
\end{tabular}

\subsection*{Experiments Retained}

\begin{itemize}
    \item Native vs.\ FUSE overhead comparison across all seven workloads
    (answers RQ1 directly)
    \item Adaptive vs.\ null prefetcher comparison (answers RQ4 directly)
    \item Cross-platform Linux ext4 vs.\ macOS APFS comparison
    \item Confidence score behavior across workload types
    (sequential $\approx 0.999$, irregular $= 0.000$, mixed $\approx 0.500$)
\end{itemize}

\section*{Revised Thesis Statement}

Traditional filesystem prefetchers control prefetching through implicit
heuristics --- detecting sequential patterns and adjusting readahead window
sizes --- but react slowly when workload behavior changes and expose no signal
an operator or policy can act on directly. JazzyFS replaces this implicit
control with an explicit, per-decision confidence value that drives prefetching
suppression immediately when confidence drops, reducing unnecessary I/O and
improving responsiveness to workload phase changes without kernel modification.
It introduces three contributions at the filesystem level: (1) sonification as
an operator-facing interface for human-in-the-loop control of prefetching
behavior, enabling a closed-loop system no existing filesystem has implemented;
(2) confidence decay rate as a metric that quantifies how quickly the prefetcher
detects a phase change, enabling faster suppression of wasted prefetching than
implicit heuristics allow; and (3) confidence trajectory shape as a
training-free workload classifier that inverts the standard use of confidence
--- rather than predicting the next access, it identifies the type of workload
currently running, enabling proactive prefetch control before failure evidence
accumulates.

\end{document}
